% Lösungen Einheit 2 von 5 – Extrempunkte (zusammenbau v0.1)
\einheitenkopf{Einheit 2 von 5 · Extrempunkte}

\erg{11}{a) Stelle gegeben – nachweisen: $x = 2$ einsetzen. \quad b) $f'(x) = 3x^2 - 3x - 6 = 0 \Leftrightarrow x^2 - x - 2 = 0$, also $x_1 = 2$; $x_2 = -1$ \quad c) $f'(x) = 3x^2 + 3x - 6 = 0$ für $x_1 = -2$, $x_2 = 1$; $f''(x) = 6x + 3$; $f''(-2) = -9 < 0$: Hochpunkt $H(-2 | 11)$; $f''(1) = 9 > 0$: Tiefpunkt $T(1 | -2{,}5)$ \quad d) $f'(x) = 4x^3 - 16x = 4x(x^2 - 4) = 0$ für $x_1 = 0$, $x_2 = -2$, $x_3 = 2$; $f''(x) = 12x^2 - 16$; $f''(0) = -16 < 0$: Hochpunkt $H(0 | 0)$; $f''(\pm 2) = 32 > 0$: Tiefpunkte $T_1(-2 | -16)$ und $T_2(2 | -16)$ \quad e) $f'(x) = 15x^4 - 60x^2 = 15x^2(x^2 - 4) = 0$ für $x_1 = 0$, $x_2 = -2$, $x_3 = 2$; $f''(x) = 60x^3 - 120x$; $f''(-2) = -240 < 0$: $H(-2 | 64)$; $f''(2) = 240 > 0$: $T(2 | -64)$; $f''(0) = 0$ und $f'''(0) = -120 \ne 0$: Sattelpunkt $S(0 | 0)$ \quad f) $f'(x) = 4x^3 - 12x^2 - 4x + 12$; $(x^3 - 3x^2 - x + 3) : (x - 3) = x^2 - 1$, also $x_1 = 3$, $x_2 = -1$, $x_3 = 1$; $f''(x) = 12x^2 - 24x - 4$; $f''(3) = 32 > 0$: $T_1(3 | -9)$; $f''(-1) = 32 > 0$: $T_2(-1 | -9)$; $f''(1) = -16 < 0$: $H(1 | 7)$ \quad g) $f'(x) = -\mathrm{e}^{x} + (3 - x)\,\mathrm{e}^{x} = (2 - x)\,\mathrm{e}^{x} = 0$; da $\mathrm{e}^{x} > 0$, ist $x = 2$; $f''(x) = (1 - x)\,\mathrm{e}^{x}$, $f''(2) < 0$; $H(2 | \mathrm{e}^{2})$, also $H(2 | 7{,}39)$ \quad h) $f'(x) = (1 - 0{,}5x)\,\mathrm{e}^{-0{,}5x} = 0$ für $x = 2$; $f(2) = 2\,\mathrm{e}^{-1}$, also $H(2 | 0{,}74)$; laut Abbildung ist es der höchste Punkt des Graphen, also ein Hochpunkt. \quad i) $r'(x) = 0{,}75x^2 - 3x + 2{,}25 = 0 \Leftrightarrow x^2 - 4x + 3 = 0$, also $x = 1$ (lokales Maximum, $r(1) = 2$) und $x = 3$ (lokales Minimum, $r(3) = 1$); Randwerte $r(0) = 1$, $r(4) = 2$. Der größte Radius $2$ cm liegt bei $x = 1$ und am Rand bei $x = 4$. \quad j) $f'(x) = 4x^3 - 4x = 4x(x^2 - 1) = 0$ für $x = 0$, $x = \pm 1$; Tiefpunkte $T(\pm 1 | -1)$, Hochpunkt $H(0 | 0)$; $f(x) \to +\infty$ für $x \to \pm\infty$; Wertebereich $W = [-1; \infty[$ \quad k) $f'(x) = 2(x - 2)\,\mathrm{e}^{-0{,}5x} - 0{,}5(x - 2)^2\,\mathrm{e}^{-0{,}5x} = (-0{,}5x^2 + 4x - 6)\,\mathrm{e}^{-0{,}5x} = 0 \Leftrightarrow x^2 - 8x + 12 = 0$, also $x_1 = 2$, $x_2 = 6$; $f''(2) = 2\,\mathrm{e}^{-1} > 0$: $T(2 | 0)$; $f''(6) = -2\,\mathrm{e}^{-3} < 0$: $H(6 | 16\,\mathrm{e}^{-3})$, also $H(6 | 0{,}80)$}
\erg{12}{a) Stelle gegeben – nachweisen: $x = -1$ in $f'$ und $f''$ einsetzen. \quad b) $f'(x) = 6x^2 - 12x$; $f'(2) = 24 - 24 = 0$ \quad c) $f'(x) = 3x^2 - 27$, $f'(-3) = 27 - 27 = 0$; $f''(x) = 6x$, $f''(-3) = -18 < 0$, also Hochpunkt. \quad d) $f'(x) = -3x^2 + 6x$, $f'(2) = 0$; $f''(x) = -6x + 6$, $f''(2) = -6 < 0$; $f(2) = -8 + 12 + 1 = 5$, also ist $H(2 | 5)$ Hochpunkt. \quad e) $f'(x) = 3x^2 - 6x = 3x(x - 2)$, $f'(0) = f'(2) = 0$; $f'(-1) = 9 > 0$, $f'(1) = -3 < 0$, $f'(3) = 9 > 0$: bei $0$ Wechsel von $+$ nach $-$ (Hochstelle), bei $2$ von $-$ nach $+$ (Tiefstelle). \quad f) $f'(x) = x^2 - 9$, $f'(3) = 9 - 9 = 0$; zusammen mit $f''(3) \ne 0$ ist $3$ eine Extremstelle. Den Funktionswert $f(3)$ auszurechnen, wäre kein Nachweis. \quad g) $f'(x) = 3x^2 - 12x + 12 = 3(x - 2)^2$, also $f'(2) = 0$; doppelte Nullstelle, $f'$ ist sonst positiv und wechselt das Vorzeichen nicht – kein Extrempunkt, sondern der Sattelpunkt $S(2 | 8)$. \quad h) $f'(x) = 4x^3 + 6x - 2$, $f''(x) = 12x^2 + 6 > 0$: $f'$ ist streng monoton steigend und hat höchstens eine Nullstelle; $f'(0) = -2 < 0$, $f'(1) = 8 > 0$, also genau eine Nullstelle mit Wechsel von $-$ nach $+$ – genau ein Tiefpunkt. \quad i) $f'(x) = 3x^2 - 5$; $f'(1) = -2 < 0$, $f'(2) = 7 > 0$; $f'$ wechselt dazwischen von $-$ nach $+$, also liegt zwischen $1$ und $2$ eine Tiefstelle. \quad j) $f'(x) = \frac{2x}{x^2 + 4} = 0$ nur für $x = 0$; für $x < 0$ ist $f'(x) < 0$, für $x > 0$ ist $f'(x) > 0$, also Tiefpunkt $T(0 | \mathrm{ln}\,4)$, also $T(0 | 1{,}39)$. Auch ohne Ableitung: $x^2 + 4$ ist bei $0$ am kleinsten und $\mathrm{ln}$ steigt.}
\erg{13}{a) $p'(x) = -2x + 6 = 0$ für $x = 3$, Scheitel $S(3 | 5{,}5)$; Abstand $d = \sqrt{(3 - 2)^2 + (5{,}5 - 6)^2} = \sqrt{1{,}25} \approx 1{,}12 > 0{,}5$. Die Differenz der $y$-Werte allein wäre nur $0{,}5$.}
\erg{14}{a) $f(-1) = -2 + 6 = 4$, also $H(-1 | 4)$; wegen der Punktsymmetrie $T(1 | -4)$; Abstand $d = \sqrt{2^2 + 8^2} = \sqrt{68} \approx 8{,}25$}
\erg{15}{a) $f'(x) = \frac{3}{8}x^2 - 1{,}5 = 0 \Leftrightarrow x^2 = 4$; $f''(x) = \frac{3}{4}x$: $H(-2 | 2)$, $T(2 | -2)$; Steigung $\frac{-2 - 2}{2 - (-2)} = -1$ und $H$ liegt auf $y = -x$: die Gerade ist $y = -x$.}
\erg{16}{a) $x^2 - 2x - 3 = 0$: $x_1 = -1$, $x_2 = 3$; Extremstelle in der Mitte der Nullstellen: $x = 1$}
\erg{17}{a) $f'(x) = 3x^2 - 6x = 0$ für $x = 0$ und $x = 2$; $f''(2) = 6 > 0$ und $f(2) = 0$: der Tiefpunkt $T(2 | 0)$ liegt auf der $x$-Achse, der Graph berührt sie dort, ohne sie zu schneiden.}
\erg{18}{a) $f'(x) = -5x\,\mathrm{e}^{-0{,}5x^2}$: für $x < 0$ steigt $f$, für $x > 0$ fällt $f$ (Graph symmetrisch zur $y$-Achse); größter Wert $f(0) = 5$; $f(x) > 0$ und $f(x) \to 0$ für $x \to \pm\infty$: Wertemenge $]0; 5]$}
\erg{19}{a) Wäre $f$ zwischen den beiden Stellen streng monoton, könnte es an ihnen nicht denselben Wert haben; also wechselt $f$ dazwischen die Richtung, und dort liegt eine Extremstelle.}
\erg{20}{a) $g_o(x) - f(x) = 0{,}25x(x - 3)^2 \ge 0$ für $x \ge 0$; $d(x) = f(x) - g_u(x) = -0{,}25x^3 + 1{,}5x^2 - 2{,}25x + 3$, $d'(x) = -0{,}75(x - 1)(x - 3) = 0$ für $x = 1$ (Minimum, $d(1) = 2$) und $x = 3$ (Maximum, $d(3) = 3$); Randwerte $d(0) = 3$, $d(4) = 2$; also $d > 0$ auf $[0; 4]$. Kleinster Abstand im Inneren bei $x = 1$, er beträgt $2$.}
\erg{21}{a) falsch – links vom Tiefpunkt fällt der Graph, also ist $f'(1{,}9) < 0$.}
\erg{22}{a) Fehler: $f''(2) > 0$ bedeutet Linkskrümmung, also Tiefpunkt. Richtig: $f(2) = 8 - 24 + 1 = -15$, Tiefpunkt $T(2 | -15)$. \quad b) Auch an einer Sattelstelle ist die Tangente waagerecht, der Graph steigt oder fällt aber auf beiden Seiten weiter. Erst ein Vorzeichenwechsel von $f'$ (oder $f''(x_0) \ne 0$) zeigt einen Extrempunkt. \quad c) Den Scheitel der Parabel bestimmen und von seiner $y$-Koordinate die Lkw-Höhe abziehen: $p'(x) = -0{,}5x + 2 = 0$ für $x = 4$; $p(4) = 5$; Abstand $5 - 4 = 1$ m.}
